% Generated draft: requires mathematical and source review.
\documentclass[11pt]{article}
\usepackage[margin=25mm]{geometry}
\usepackage{fontspec}
\setmainfont{Latin Modern Roman}
\usepackage{amsmath,amssymb,mathtools,unicode-math}
\setmathfont{Latin Modern Math}
\usepackage{enumitem,needspace,xurl,hyperref}
\setlength{\parindent}{0pt}
\setlength{\parskip}{5pt}
\setlength{\emergencystretch}{3em}
\newcommand{\N}{\mathbb{N}}
\newcommand{\Z}{\mathbb{Z}}
\newcommand{\Q}{\mathbb{Q}}
\newcommand{\R}{\mathbb{R}}
\newcommand{\C}{\mathbb{C}}
\newcommand{\F}{\mathbb{F}}
\newcommand{\A}{\mathbb{A}}
\newcommand{\PP}{\mathbb{P}}
\newcommand{\E}{\mathbb{E}}
\newcommand{\eps}{\varepsilon}
\newcommand{\dd}{\,\mathrm{d}}
\begin{document}
\section*{UnsolvedMath ID: 30003226}
\textbf{Source problem code:} catalogue entry
\subsection{Definitions and mathematical statement}
The supplied record asks whether poles of the motivic zeta function of a polynomial are roots of its Bernstein--Sato polynomial. The coefficient field, the version of the motivic zeta function, the meaning of a motivic pole, and the local or global scope are not specified.

For reference, in the conventional complex setting, let \(n\geq 1\) and let \(f\in\C[x_1,\ldots,x_n]\) be nonzero and nonconstant. Write \(D_{\A^n_{\C}}\) for the ring of polynomial differential operators. The global Bernstein--Sato polynomial \(b_f(s)\) is the monic generator of the nonzero ideal
\[
\left\{b(s)\in\C[s]: P(s)f^{s+1}=b(s)f^s\text{ for some }P(s)\in D_{\A^n_{\C}}[s]\right\}.
\]
The identity is interpreted in the formal module \(\C[x_1,\ldots,x_n,f^{-1}][s]f^s\), with
\[
\partial_{x_i}(g f^s)=\left(\partial_{x_i}g+s g f^{-1}\partial_{x_i}f\right)f^s.
\]
A local Bernstein--Sato polynomial instead uses differential operators with coefficients in a specified local ring at a specified point; neither that point nor a local convention is supplied.

Motivic zeta functions involve the class \(\mathbb L\) of the affine line in an appropriate Grothendieck ring, often localized or completed. With the formal substitution \(T=\mathbb L^{-s}\), a denominator factor
\[
1-\mathbb L^{-\nu}T^N,\qquad N>0,
\]
corresponds to the candidate value \(s=-\nu/N\). Such a denominator factor alone does not establish a genuine pole. A precise pole convention and coefficient ring must be fixed; motivic expressions cannot simply be treated as ordinary complex meromorphic functions.

\textbf{Proposed implication, with unresolved conventions.} For every polynomial in the intended class, does every genuine motivic pole \(\alpha\), expressed in the \(s\)-parameter, satisfy
\[
b_f(\alpha)=0?
\]
Under the complex global interpretation above, this would ask, for every \(n\geq1\) and every nonzero nonconstant \(f\in\C[x_1,\ldots,x_n]\), whether
\[
\alpha\in\operatorname{Poles}_s(Z_f^{\mathrm{mot}})\quad\Longrightarrow\quad b_f(\alpha)=0.
\]
Here \(Z_f^{\mathrm{mot}}\) and its pole set still require the missing choices. The complex global interpretation and the exclusions of zero and constant polynomials are explicitly conditional conventions, not hypotheses supplied by the record. The intended conclusion concerns actual Bernstein--Sato roots, rather than merely their exponentials as monodromy eigenvalues.
\subsection{Short English statement}
The question asks whether every genuine pole of a polynomial's motivic zeta function is a root of its Bernstein--Sato polynomial. The source leaves the field, local or global scope, zeta function variant, and pole convention unspecified, so a unique fully precise formulation requires review. Its supplied open status is a dated curation claim.
\subsection{Sources}
\begin{itemize}
\item Supplied local UnsolvedMath registry record, permanent ID \(30003226\), source problem code \texttt{catalogue entry}; registry paths identified by the project are \texttt{data/unsolvedmath/problems.json} and the supplementary \texttt{catalog/problems.json}. The record belongs to the Oberwolfach Reports Open Problems collection and supplies proposed year 2016; no individual proposer is named.
\item Supplied report citation: \emph{Singularities (2017)}, Oberwolfach Reports, pp.\ 2625--2691, DOI \texttt{10.4171/owr/2016/46}. \url{https://doi.org/10.4171/owr/2016/46}. The same link is supplied in the literature list as \emph{Oberwolfach Report 46}. The citation's 2017 designation and the record's 2016 year information are retained without reconciliation.
\item Supplied literature reference: \emph{On the monodromy conjecture for non-degenerate hypersurfaces}. \url{https://doi.org/10.1007/s00209-019-02342-5}.
\item Literature assessment dated exactly \texttt{2026-08-21}: the supplied curation labels the problem open, describes the extracted claim as the general Igusa--Denef--Loeser monodromy conjecture, reports substantial special-case results, and says the general polynomial statement remains open, including unresolved surface cases. This is supplied dated curation derived from the OWR material, not an independently verified literature assessment. Its terminology and scope require checking against the cited sources.
\item The registry supplies statement status \texttt{corrected\_verified} and reports a manual semantic, typographic, and LaTeX review of the clean statement. These are supplied registry annotations, not independent verification of this draft or resolution of its formulation issues.
\end{itemize}
\end{document}
